\clearpage\thispagestyle{empty}
\pdfbookmark[0]{Sources for Comet (18 February 1819)}{cometfeb1819-sources}
\begin{center}\textbf{Sources for Comet (18 February 1819)}\end{center}
\noindent This AI translation follows \textit{Comet} in \href{https://gdz.sub.uni-goettingen.de/id/PPN236007467}{\textit{Werke}, vol.~VI} (Göttingen, 1874), pp.~417--419. The report first appeared under the heading Göttingen in \textit{Göttingische gelehrte Anzeigen}, no.~28, 18 February 1819, pp.~\href{https://books.google.com/books?id=CAQEAAAAMAAJ&pg=PA273}{273}--278. The University of Michigan copy digitized by Google was consulted for pp.~273--\href{https://books.google.com/books?id=CAQEAAAAMAAJ&pg=PA277}{277}. The report refers to Gauss in the third person.
\par\smallskip\noindent The first printing gives 9,52933 for the logarithm of the least distance in Nicolai’s provisional parabolic elements on p.~\href{https://books.google.com/books?id=CAQEAAAAMAAJ&pg=PA274}{274}, where the collected edition has 0.52933. The date 1818 November and the wording “longitude of the ascending comet” in the comparison with the comet of 1805 also occur on p.~\href{https://books.google.com/books?id=CAQEAAAAMAAJ&pg=PA276}{276} of the first printing. The translation follows the collected readings, including the figure “about 40” without a unit in the right-ascension comparison. The historical footnote about probable Seeberg time also appears on p.~\href{https://books.google.com/books?id=CAQEAAAAMAAJ&pg=PA275}{275} of the first printing. Translation and typesetting produced by AI.
\clearpage
